% Chapter 10: Cézanne's Space; Rigid Perceptual Perspective
% Source pages: 215–232

\chapter{C\'ezanne's Space; Rigid Perceptual Perspective}

\markboth{Chapter X. C\'ezanne's Space}{Chapter X}

The work of Cézanne invariably attracts the attention of artists and art historians.
This master entered the history of world painting as an artist who blazed new trails,
boldly breaking familiar and seemingly inviolable rules and conceptions. If one
restricts oneself here to the single question of Cézanne's perspective system, the
consensus today on this side of his art amounts to the view that he destroyed the
system of perspective created in the Renaissance and for more than four hundred years
faithfully serving European painting. Fritz Novotny gave his interesting work on
Cézanne the very title ``Cézanne and the End of Scientific
Perspective''.\footnote{F.~Novotny. \emph{C\'ezanne und das Ende der
  wissenschaftlichen Perspektive}. Wien, 1938.}

The bare fact of the ``destruction'' of the system of linear perspective only shows
that for some reason this ``academic'' (by the 19th century) system did not satisfy
Cézanne. The writings of art historians contain many varied and often persuasive
thoughts and observations connected with the desire to understand the geometric
peculiarities of Cézanne's spatial constructions. Many who have written about this
artist agree that in departing from linear perspective Cézanne thereby departed from
formal-geometrically ``correct'' depiction and moved toward depictions that, while
artistically justified, no longer rest on a strict geometric foundation. On the other
hand, a number of Cézanne's own statements and the recollections of his
contemporaries suggest that he strove to transmit nature as accurately as possible.
``Truth is in nature; I will prove it,'' he
said.\footnote{A.~Perruchot. \emph{C\'ezanne}. Moscow, 1966, p.\,323.}
Attempting to explain his puzzling perspective constructions, Cézanne as a rule
argued them by saying that he saw things that way. This in turn gave rise to the
suggestion that he suffered from a visual defect. Probably not wishing to waste time
on idle argument, Cézanne did not deny such a suggestion.

\srcpage{215}

If today one returns to the disputes of the early century about Cézanne's perspective
``errors'', armed with the results obtained in previous chapters, one can advance the
following conjecture. As was shown above, the system of linear perspective gives a
distorted depiction of space---man sees by far not according to its rules. Endowed
with an especially acute visual perception of space, which the 19th century
perspective school was unable to suppress, the artist could not but feel this---he
transitioned to the direct transmission of visual perception on his canvases. This
means that Cézanne began working in the system of perceptual perspective. But then
his perspective constructions have strict geometric foundation, no less strict than
the system of linear perspective. Cézanne may not have known this scientific system,
but it is remarkable in that it can be followed intuitively, for the single thing it
demands of the artist is to follow his visual perception as accurately as possible.

Let us turn to an analysis of the geometry of Cézanne's paintings and compare it
with the same subjects executed in strict linear perspective. Such a comparison is
enormously facilitated by the already-mentioned monograph of Novotny, which
contains not only reproductions of Cézanne's paintings but photographs of the
sections of landscape he depicted, taken from the same viewpoints from which he
painted them.

\srcpage{216}

Before proceeding to the programme outlined here, however, it is necessary to take
one further step in the theory of perceptual perspective systems. In studying
medieval art, the problem of an adequate-to-visual-perception depiction of space
from foreground to horizon, over a sufficiently large angle of view, was in
substance never raised; and without consideration of this question an analysis of
the geometry of Cézanne's paintings is simply impossible.

It is appropriate to recall here that in this book the geometric structure of some
method of depicting space on a plane is constructed analytically, prior to any
engagement with works of art, and then an attempt is made to identify the features
so found in the works of a given master. This approach has the advantage of high
objectivity, since it completely excludes the possibility of even unconsciously
``imposing'' on the master impulses alien to him but close to the person studying
his work.

In the preceding chapters, when speaking of the system of perceptual perspective,
the discussion mainly touched only on the question of depicting individual objects.
This was entirely sufficient for the analysis of Byzantine and Old Russian painting,
as well as of Iranian and Indian miniatures, whose geometry can as a rule be treated
as a developed system of local axonometries and transformations of those axonometric
structures. Let us now pose the task of creating a perceptual perspective system
that would be subject to precisely formulated mathematical rules and would possess
the property characteristic of the system of linear perspective---that the person
using it (not necessarily an artist!) could depict any pictorial space on paper
according to completely clear, unambiguous rules.

At first glance the task posed contradicts the basic property of the system of
perceptual perspective discussed so many times above: the right of the artist to
freely choose distortions. Let us first confirm here too that depicting space
without distortions is fundamentally impossible, and consequently the discussion
concerns only the manner---now we shall speak of a pre-specified manner---in which
these inevitable distortions are to be chosen.

\srcpage{217}

In medieval art the artist, depicting an individual object, each time freshly
resolved this task, deciding which particular geometric feature of that object he
would show ``incorrectly''. If he unthinkingly followed tradition, this naturally
does not change matters, for here the particular artist merely repeated an already
established and confirmed device. In this case, the ``distortions'' within a single
work could be determined to some degree by considerations that were mutually
independent of one another (speaking only of geometric constructions).

If one agrees in advance to choose the character of distortions, the situation may
be entirely different. Let us pose the task (as in the theory of linear perspective)
of depicting an individual point of the space of objects. Then the inevitable
distortions can manifest only in the depiction of the distances of that point from
the picture plane (depiction of depth), from the object plane (depiction of height),
and from the vertical plane normal to the picture plane and containing the principal
point of the picture (depiction of the lateral position of the point). If one
specifies precisely not only which of these three segments is to be shown with
distortion, but also states exactly (mathematically, that is, numerically!) how and
in which direction this distortion is to be made, then a perceptual perspective
system can thereby be created in which no room remains for the artist's
``arbitrariness''. Such a system is appropriately called rigid, as was done in
Chapter~III.

A rigid perceptual perspective system will differ from a free one in that the
artist's freedom is not limited before the work begins---while he decides what can
most painlessly be distorted and in what manner. Once the corresponding decision has
been freely made, the strict law of geometric construction resting on mathematical
rules comes into force. It is of course clear that a rigid perceptual perspective
system---unlike the system of linear perspective---is not uniquely determined. In
principle one can shift the inevitable distortions onto any one of the three segments
determining the position of a point in pictorial space, or onto a combination of
them.

Of course, the choice of distortions should be made so that they are as imperceptible
as possible and affect as few elements of the image as possible. An example of one
such choice, with a clear demonstration of what and how to ``distort'', is given in
Appendix~8. There it is also shown how the attempt to render everything without
distortion inevitably leads to ruptures in the image, since then the depiction of
any point of pictorial space ceases to be unique and proves to lie simultaneously at
different points of the picture plane. Thus the ``primacy of continuity'' mentioned
in previous chapters has a deep mathematical justification, whose discussion is out
of place here.

As an example of one of the possible rigid perceptual perspective systems, let us
consider here in more detail the one in which the inevitable distortions are
connected with the depiction of heights; this is precisely the one discussed in the
aforementioned Appendix. It is obvious that such a system, correctly\footnote{Here
  and below, the word ``correctly'' means: in accordance with visual perception.}
transmitting both horizontal dimensions, will be especially suited to the depiction
of landscapes in which verticals play a subordinate role.

Here one should, however, specify that all these distortions concern only the
foreground and middle distance, since, as has been noted many times, perceptual
perspective can render distant planes without distortion. Therefore this
``landscape'' variant of the rigid perspective system is able to render without
distortion not only a flat landscape but also a landscape with mountains in the far
distance.

One should not think that once the decision has been taken to shift the inevitable
distortions onto the depiction of heights, everything else will be rendered without
distortion. Even with a correct depiction of horizontal dimensions, the question of
the correctness of the depiction of angles remains open. Here the problem is in some
sense hopeless, for in Chapter~III it was already shown that a perspective system
correctly rendering all three angles at a vertex of a parallelepiped cannot exist.
Consequently, as in the system of linear perspective, in other possible perceptual
perspective systems, and in the system introduced here, angles will be rendered
partially correctly and partially not. Probably this circumstance does not play an
especially large role; relative sizes of objects in a picture attract attention
rather more than inaccuracies in foreshortening. It is not excluded that this has
deeper roots in the human psyche: peoples at the most primitive stages of development
already use measures of length but know nothing of angular measures, which appear
comparatively late.

\srcpage{218}

The considerations set out above allow one to classify various rigid perspective
systems on the basis of only the correctness of the depiction of lengths of straight
segments. To give a vivid comparative evaluation of different perspective systems,
we shall use the following device. Let us divide all the space to be depicted into
three planes---near, middle, and far (to the horizon)---and examine what the degree
of correctness of the depiction of a cube of small size at these three distances
will be. Let the cube be positioned so that its faces are parallel to the three
mutually perpendicular natural planes of pictorial space (ill.\,98). In this case
the edges $a$, $b$, and $d$ will give a clear representation of the three principal
directions: laterally, upward, and in depth from some initial point~$O$ lying in the
object plane. The apparent size of edge~$d$ depends on the lateral displacement of
the cube---this follows from comparing the depictions~I and~II in ill.\,98. It is
easy to see that the apparent distance between the front and rear faces of cubes~$d$
is the same for both cubes. Consequently, it is reasonable to transition from the
indeterminate~$c$ to the constant~$d$ for all equally distant cubes. If alongside
$a$, $b$, $d$ one also considers $a'$, $b'$, $d'$, this will allow one to judge the
geometry of the depiction as a whole. Numerical coincidence or difference of the
values $a$, $b$, $d$, $a'$, $b'$, $d'$ with the corresponding values computed from
the formulas describing visual perception indicates which elements of the depiction
are rendered correctly and which are (and to what degree) distorted.

We shall evaluate the degree of correctness of the depiction according to the
following nine quantities: $a$, $b$, $d$, $a/b$, $b/d$, $d/a$, $a'/a$, $b'/b$,
$d'/d$. The first three quantities will indicate how correctly (in terms of length)
the corresponding segment is depicted, while the remaining quantities will indicate
how correctly the similarity of the three faces of the cube is preserved in the
depiction (independently of scale correctness).

\figblock{98}{Basic dimensions used for the numerical evaluation of the correctness
  of depiction. (A cube on the near, middle, and far planes, with edges $a$
  (lateral), $b$ (vertical), $d$ (depth).)}

Let us compare the system of linear perspective and the variant of the rigid
perceptual perspective system taken above as an example. Such a comparison is
legitimate since the system of linear perspective is also rigid. To carry out the
comparison in practice, one must agree on what size will be taken as the reference
standard of ``correct depiction''. If for clarity one turns to photography, it is
obvious that from a single negative one can print both a small and a greatly
enlarged photograph, and both will have the same geometry, although at different
scales. Therefore one should agree to evaluate all sizes in the photograph relative,
say, to the length of the longer side of the frame; numerical comparison of two
photographs at different scales then becomes possible. We shall analogously take as
such a reference standard (module) the value~$a$ at the middle plane, and judge the
correctness of other segments relative to it, as if bringing the images in the
drawings to a single scale.

\srcpage{219}

The results of the comparison of the two perspective systems are set out in the
table. In it the sign ``$+$'' means correct depiction and ``$-$'' means
distorted.\footnote{In drawing up the table, 9 parameters were chosen as
  determining ones for each plane. This was motivated by the desire for the most
  vivid presentation of the results; a mathematical analysis of the independence of
  the parameters introduced was not conducted. Although linear perspective can be
  regarded as a special case of perceptual perspective, it is here deliberately
  separated from it, since the main point is the comparison of the traditional
  perspective constructions of the 19th century with the constructions introduced by
  Cézanne.}

\begin{center}
\fbox{\parbox{0.88\textwidth}{\smallskip\centering
\textbf{Table.} Comparison of characteristics of linear perspective and the rigid
perceptual perspective system.\\[0.5ex]
{\small[Rows: Linear perspective; Perceptual perspective.
Columns: 9 geometric parameters of a cube ($a$, $b$, $d$, $a/b$, $b/d$, $d/a$,
$a'/a$, $b'/b$, $d'/d$) $\times$ 3 planes (near, middle, far) $=$ 27 columns.
``$+$''~$=$~correct depiction; ``$-$''~$=$~distorted.
Linear perspective: 15 distortions; perceptual perspective: 8 distortions.
Source: p.\,219.]}
\smallskip}}
\end{center}

The data in the table show that out of 27 parameters, 15 are distorted in the system
of linear perspective---more than half---while in the perceptual system only 8 are
distorted---less than one third. Consequently, the latter system will be
significantly closer to natural visual perception. In essence the sole cause of
distortions in the system of perceptual perspective (let us recall: inevitable) is
the distortion, in accordance with the choice made, of the sizes of vertical segments
(parameter~$b$) on the near and middle planes. As for the system of linear
perspective, the distortions inherent to it are significantly more substantial.
First, it does not give the correct scale of depictions on the near and far planes
(parameters $a$, $b$, $d$) and, second, it distorts the configurations of the cube
faces on the near and middle planes more substantially than the perceptual system
does (4 vs.\ 3 distorted parameters). One can only marvel at how many distortions
of natural visual perception are carried by the perspective system that for centuries
was considered the ideal of accurate, scientific representation of space on the
picture plane! (The mathematical justification of the table is given at the end of
Appendix~8.)

To get a clearer sense of the character of the distortions inherent to both
perspective systems, let us compare the most fundamental: the non-preservation of
the relative scales of depictions. Let us cite some purely approximate numerical
estimates, using the example considered in Appendix~8. In the system of linear
perspective the distortion of scales of the near and far planes relative to the
middle amounts to roughly 35\%, and must therefore be considered substantial. In the
system of perceptual perspective, the distortion of scales of vertical segments (in
the direction of their enlargement---parameter~$b$) amounts to roughly 20\% for the
beginning of the middle plane and 40\% for the near plane, while no distortion of
horizontal segments occurs.\footnote{Here and below the natural assumption is made
  that the near plane of landscape painting lies beyond the axonometric zone of
  close space; the latter is therefore excluded from consideration.}
However, if there are no tall objects on the near plane, the distortion of scales of
vertical segments on the near plane will not manifest itself, and the picture will
on the whole prove much closer to natural visual perception than one executed in the
system of linear perspective---whose distortions affect not only vertical but also
horizontal segments. In addition, one should not forget that in the system of linear
perspective the distortions of object configurations are also more substantial.

\srcpage{220}

The general comparison of the two rigid perspective systems---linear and
perceptual---is usefully concluded by comparing images built in each of the two
systems. Figs.\,99 and~100 give an image of a conventional landscape: the first
rendered in linear perspective, the second in perceptual perspective. The middle
plane in both images (line~$AA$) is shown at the same scale to facilitate comparison.

\figblock{99}{Conventional landscape rendered in the system of linear perspective.}

\figblock{100}{Conventional landscape rendered in the system of perceptual
  perspective.}

\srcpage{221}

Let us note first what the two images have in common. Both convey depth of space,
volumetric modelling of form (though not identically well), and both images depict
objects and the intervals between objects according to identical rules. Both images
give examples of complete perspective systems with their unifying spatial approach
to the rendering of depth on the picture plane, with a single viewpoint in both
perspective systems. Perhaps it is not superfluous to recall here that traditionally
all of these properties were attributed exclusively to the system of linear
perspective.

To characterise briefly the difference between the two perspective systems, one can
assert---taking the familiar system of linear perspective as the baseline---that in
the perceptual perspective system the sizes of objects in the far plane (the
mountains) are enlarged relative to the baseline, and those on the near plane (the
flower bed) are reduced. The depth of the middle plane (from line~$AA$ to the
mountains) has increased compared to the baseline (the area of the bay shown on the
right has therefore also increased), while the depth of the foreground (from
line~$AA$ to the lower edge of the picture) has decreased. In addition, objectively
straight lines (which are accordingly straight in the image built by the laws of
linear perspective) have become curvilinear (the road leading to the mountains, the
hedge depicted behind the house).

These are precisely the features characteristic of Cézanne.

Having examined Cézanne's landscapes and corresponding photographs of the depicted
locations, Novotny summarises the results of such a comparison in the following
words: ``In many cases changes in relative sizes are made, such that perspectively
large foreground objects are rendered somewhat reduced and, conversely, far-distance
objects are exaggerated in the
picture.''\footnote{F.~Novotny. Op.~cit., S.\,39.
  An analysis of Cézanne's canvases and the corresponding photographs included in
  Novotny's monograph---measuring comparable objects in the near, middle, and far
  planes---showed that Cézanne enlarges far-plane objects by an average of 25\%
  (relative to photographs) and reduces near-plane objects by 25--40\% relative to
  middle-plane sizes. This is in good numerical agreement with the corresponding
  sizes obtained in ills.\,99, 100.}

\srcpage{222}

A vivid impression of this characteristic of Cézanne's landscape painting can be
gained from his canvas ``The Aqueduct'' (ill.\,101). The Mont Sainte-Victoire shown
near the horizon reminds us---accustomed as we are to photography---of the
Himalayas rather than a modest elevation of France, while the foreground area is
clearly ``compressed''. As for the trees, they are (as is to be expected in the
landscape variant of the rigid perspective system under consideration) shown at
exaggerated rather than normal height.

This mode of landscape depiction produces a kind of surprise in Novotny. Noting
that Cézanne constantly strives for a ``portrait'' precision in rendering the
landscape, he attempts to connect these violations of the laws of linear perspective
with their being made ``in favour of colour''.\footnote{F.~Novotny. Op.~cit.,
  S.\,5.} However, the artist's ``reduction of the intensity of impact of the most
elementary means of conveying spatiality''---a kind of softening of the spatial
accent\footnote{Ibid., S.\,4.}---appears to him an artistic device, though one
without rational explanation.

\figblock{101}{C\'ezanne. ``The Aqueduct''. 1887--1890. Pushkin Museum of Fine Arts,
  Moscow.}

\figblock{102}{Scheme of the depiction of an ``avenue'' in the systems of linear
  and perceptual perspective.}

In the light of the comparison of images of the conventional landscape in the two
perspective systems made above, these ``peculiarities'' of Cézanne become entirely
comprehensible. It is precisely the striving for ``portrait'' precision in his
landscapes, the strict adherence to visual perception (that is, the use of the
perceptual perspective system), that led to the deformations of the familiar
perspective that puzzled art historians. These deformations had as their aim not
moving away from the regularities of natural visual perception but, on the contrary,
moving closer to them.

That the shift to a perceptual perspective system has taken place is indicated not
only by the general character of the perspective deformations, but also by details
of the depiction. Thus, the ``reduction of convergence'' noted by Novotny in the
depiction of the edges of an avenue shown receding from the
viewer\footnote{Ibid., S.\,37, Tabl.\,37.} is nothing other than the natural
depiction of such an avenue in the perceptual perspective system. The essence of
this transformation is clear from schematic ill.\,102, where lines $B'O$ give the
depiction of the avenue---having width $AA$ at the middle plane---that it would have
in the system of linear perspective, while lines $AB$ convey the configuration of
such an avenue in Cézanne. It is quite clear that the ``unfolding'' of the avenue
edges from position $B'O$ to $AB$ and consequently the foreshortening of the near
part of the avenue from dimension $B'B'$ to $BB$ has the same origin as the
analogous changes at middle and near planes in the transition from the left image to
the right in ill.\,17 (from linear to perceptual perspective). Thus here too one
should speak not of moving away (as Novotny supposes) but of bringing the geometry
of space shown in the picture closer to natural visual perception.

Speaking of the peculiarities of the depiction of the foreground in Cézanne's
landscapes, Novotny (and other researchers) particularly note that this part of the
depicted space as it were sinks downward, ``falls away'', which together with the
contraction of its dimensions in the picture gives a special originality to the
artist's works. But such ``sinking'' is a natural consequence of the recourse to
the perceptual perspective system: its causes were discussed in detail in
Chapter~III, where it was then illustrated by using as an example the painting by
V.~Bednov, ``Bleaching Canvases'' (ill.\,21). Thus this feature of Cézanne's work
too becomes comprehensible and natural within the framework of the theory of
perceptual perspective.

\srcpage{224}

\figblock{103}{C\'ezanne. ``Chestnut Avenue at Jas de Bouffan''. Watercolour.
  1883--1887. Paris. Former collection of Paul Cézanne fils.}

Some peculiarities in the perception of Cézanne's paintings are connected with the
fact that they, having been painted in the system of perceptual perspective, are
perceived by viewers trained in the system of linear perspective. As an example let
us consider the question of the artist's elevated viewpoints. In many studies
devoted to Cézanne's work, his striving to paint from ``unnatural'', elevated
viewpoints is particularly emphasised. The arising of this impression will be
explained by returning to the discussion of the depiction of the avenue (ill.\,102).
Let the actual height of the artist's viewpoint be represented by vanishing point~$O$.
As a result of the fact that in the perceptual perspective system the edges of the
near end of the avenue $B'B'$ shift to $BB$, a new vanishing point $O'$ for the
avenue edges arises. A person accustomed to the system of linear perspective will
perceive a picture in which the avenue width $BB$ is shown as painted from an
unnaturally high viewpoint corresponding to $O'$; in such a viewer's perception the
artist seems to have ``taken flight'' above the ground by the amount $OO'$ in order
to paint his landscape. Yet this is not the case at all. For the vanishing point~$O'$
is obtained by extending the edges of the avenue $AB$ as straight lines, whereas in
the perceptual perspective system these edges can only be considered rectilinear
approximately, over limited stretches, and when extended from the foreground to the
horizon must be depicted as curves, analogous to what is shown in ill.\,17. But
then $O'$ is not at all a point lying on the horizon, and consequently does not at
all characterise the actual height of the artist's viewpoint.

\srcpage{225}

It may turn out that the avenue edges $AB$, when depicted as curvilinear arcs
(shown dashed), give the same vanishing point~$O$ as in the linear perspective
system. Therefore it is more correct to say not that Cézanne painted his avenue from
an elevated viewpoint, but that it, having been painted from a normal viewpoint, is
perceived distortedly by the majority of viewers.

These general arguments can be illustrated by analysing one of Cézanne's depictions
of the avenue at Jas de Bouffan (ill.\,103). If one constructs the usual vanishing
point and draws a horizontal line through it (the horizon line), it intersects the
depictions of the nearest trees at a level corresponding, within the theory of linear
perspective, to the height of the artist's eyes. It is perfectly obvious that in
ill.\,103 this height comes out greater than a normal person's stature. If one
proceeds from the conviction that Cézanne worked in the system of perceptual
perspective, this position of the formal vanishing point can be explained by the
already-discussed scheme shown in ill.\,102, or by the fact that, taking the
foreground scale as the basis, the artist was compelled to raise the horizon line,
as follows from ill.\,17.

Thus if a painting that employs the perceptual perspective system is interpreted from
the standpoint of the theory (or the experience of perception) of linear perspective,
this inevitably leads---and does lead---to distorted conceptions. Of course, the
considerations presented here do not at all mean that Cézanne never employed elevated
viewpoints. Investigation of this question is the business of art historians; however,
in such an investigation the effect described here must also be taken into account,
distinguishing cases where the artist genuinely uses a high viewpoint from those
where it only seems so owing to the habit of the linear perspective system.

\srcpage{226}

In the light of what has been said, some of Novotny's general conclusions also
acquire a different nuance. In particular, his observation that Cézanne is
characterised by a ``reduction of intensity'' in the methods of conveying spatiality
developed by the creators of the linear perspective system can now be interpreted in
the opposite sense. One can assert that Cézanne was transmitting a completely
natural perception of space, while before him artists had exaggerated the rendering
of spatiality, inordinately enlarging objects in the foreground and reducing them in
the background. Viewers accustomed to images with such perspective exaggerations
perceive the transition to naturalness as softening (to use Novotny's phrase), that
is, as distortion of the ``correct'' spatial accent.

The question of the spheroidality of Cézanne's space deserves special consideration.
Liliane Guerry, who proposed this term, writes: ``The vision of space in
Cézanne\ldots is a spheroid mobile field. Space seems to rotate; this impression
arises from the curved lines forming the basis of the
composition.''\footnote{Cited from the article: V.\,N.\,Prokofiev. The exploit of
  the great stubborn one.---In: A.\,Perruchot. \emph{C\'ezanne}, p.\,351.}
Some impression of this spheroidality is given by the well-known Cézanne canvas
``The Road Turning'' (ill.\,104). From the quotation cited it follows that the
``curved lines'' are a compositional device of the artist and consequently a question
that does not form the subject of study in the present book. Therefore below the
problem of the spheroidality of Cézanne's space will be considered in a narrow sense,
only to clarify how natural these ``curved lines'' are in the system of perceptual
perspective and whether they have, at least in part, a rational origin.

\figblock{104}{C\'ezanne. ``The Road Turning''. 1879--1882. Museum of Fine Arts,
  Boston.}

\srcpage{227}

\figblock{105}{Scheme of the curvature of the zero line.}

First of all one should recall what was said above in comparing the conventional
landscapes in ills.\,99 and~100. In the system of perceptual perspective, objectively
straight lines are depicted in the form of curves. This circumstance was already
discussed in Chapter~III. Consequently, the very fact of the curving of objectively
straight lines in the system of perceptual perspective is entirely natural. It is
equally natural that such curvature endows the depicted space with a certain
sphericity. As follows from ill.\,100, a series of lines whose convexity is directed
upward arises in the picture; this can be interpreted as rendering the subconsciously
sensed convexity of the earth's surface---although formally and mathematically it is
the inevitable consequence of reconciling the scales of horizontal segments in the
near, middle, and far planes, while simultaneously requiring that far planes be
rendered in full accordance with visual perception (since this is in principle always
possible).

Let us continue the analysis.

When the artist, for the purpose of refining his visual impressions, shifts the
direction of gaze from one area of space to another, the entire complex system of
retinal image processing ``moves'' together with the zone of clear vision. Above, in
Chapter~V, the concept of the zero line was introduced---the line that the artist,
for a given direction of gaze, consciously or unconsciously takes as a certain
``reference'' line, relative to which the entire adjacent space to left and right is
constructed. Let this zero line be an objectively straight line, seen in
foreshortening and therefore an inclined straight in the system of linear perspective.
In ill.\,105, such a zero line is shown in the perceptual perspective system as arc~$AB$
with convexity directed downward, constructed for the same conditions as the
conventional landscape in ill.\,100. In its construction the transformations of
linear quantities occurring during visual perception were taken into account in the
appropriate manner, and it correctly conveys visual perception in its main features.
Although in such a construction the linear quantities are transformed appropriately,
the angles---in particular the angles of intersection of arc~$AB$ with the bottom
edge of the picture and the horizon---are shown distortedly; the correct angles are
given by the directions $AA'$ and $BB'$.

Here again appears the well-known ``antagonism'' between the rendering of lines and
angles, of which mention has already been made more than once. Whatever the objective
causes of this ``antagonism'', the artist faces a difficult task: if he correctly
renders the course of the zero line~$AB$, he must show it with convexity downward;
if on the other hand he considers the angles of its meeting with the horizon and the
bottom edge of the picture more important, then, connecting directions $AA'$ and
$BB'$ (shown dashed), he must show it with convexity upward. If he retains the solid
line~$AB$ as the image but wishes to correctly render its meeting with the horizon,
he must ``tilt'' the horizon, giving it the direction $CC'$ at least in the
immediate neighbourhood of point~$B$.

\srcpage{228}

Since with repeated transfers of the direction of gaze and the depiction of many
zero lines (which in reality may be various extended features---road segments,
boundaries between different types of vegetation, river banks, and so on) these
will in different places and in different ways approach the horizon, the latter (as
a result of phenomena of the kind described above) may prove to be depicted convex
or concave, while the zero lines themselves ``curve'' in various ways. On a more
detailed examination of these natural ``curvings'' of objectively straight lines, it
can be shown that they are maximally curved in the middle-plane region.

To bring some order into this apparent chaos, one can paint the picture by assuming
the direction of gaze to be fixed (as is done in the system of linear perspective)
and introducing on this basis a single zero line for the picture. This was done in
ill.\,100, where the vertical passing through the principal point of the picture was
chosen as the zero line. In this way a formally ordered image is easily obtained.
However, the artist, even having constructed such a formally ordered space for the
picture as a whole, is entitled to deviate from it in depicting individual elements
of the picture, in an attempt to convey more precisely through additional curvings
one or another local feature of the landscape he is reproducing. Moreover, the
artist is of course by no means obliged to paint the picture from a single viewpoint.

Since even the striving for ``portrait'' precision in a landscape leads to a complex
and ambiguous system of curvings of objectively straight lines in the picture, it
must be recognised as entirely natural if the artist reinforces such effects for
compositional reasons. Consequently, both the spheroidality of Cézanne's space and
the diversity of its manifestations have not only purely compositional but also
objective causes. (The mathematical justification of the assertions made above can
be found in Appendix~9.)

\srcpage{229}

The influence---noted above in connection with the problem of the spheroidality of
Cézanne's space---of the depiction of local structures ``examined'' separately on
the composition of the picture as a whole is also observed in his still lifes. A
discussion of the connected range of questions can be found in Aaron
Berkman;\footnote{A.~Berkman. \emph{Art and Space}. N.\,Y., 1949, pp.\,115--123.}
however, if one examines the question of spatial constructions in Cézanne's still
lifes in an integral manner, trying to understand their connection with natural
visual perception, the main difficulty of such an analysis consists in the absence
of material for comparison (photographs of the same groups of objects as depicted in
the still lifes). Perhaps the only thing that speaks quite clearly here too of the
artist's striving to rely on the natural perception of space is the frequently
observed ``elevated viewpoint''---the apparent ``raising'' of the distant parts of
the table on which the depicted objects are arranged. This peculiarity is perfectly
comprehensible from the standpoint of the theory of perceptual perspective for near
regions of space, and is explained by the action of the form-constancy mechanism.
This can be verified by turning to the still life ``Peaches and Pears'' (ill.\,106).

\figblock{106}{C\'ezanne. ``Peaches and Pears''. 1888--1890. Pushkin Museum of Fine
  Arts, Moscow.}

\srcpage{230}

The adherence to the system of perceptual perspective is visible not only from the
``raising'' of the distant part of the table---if the artist had depicted its legs,
the far ones would have been larger than the near ones---but also from the fact that
the table surface is shown in axonometry: its sides are parallel. The mode of
depiction used corresponds precisely to the scheme in ill.\,40B.

Thus, turning to Cézanne's still lifes also does not change the overall conclusion
that in his paintings he mainly followed the regularities of the natural perception
of space---that is, the system of perceptual perspective.

The rigid perceptual perspective system developed above in the analysis of Cézanne's
landscapes (described and justified in detail in Appendix~8)---by means of which not
only an artist but any draughtsman can give a depiction of space closer to visual
perception than the system of linear perspective allows---is of course not intended
for that purpose alone. It should help the artist more deeply understand the
possibilities of rendering space on a plane and remove the ``hypnotic'' influence
of the centuries-old system of linear perspective. The new perspective system will
probably also prove useful to art historians, who will now be in a position to
evaluate the spatial constructions of artists with more justification from the
standpoint of their correspondence to visual perception.

Only one variant of the rigid perceptual perspective system was used here. It does
not cover the very close space characteristic of medieval art, whose modes of
depiction are described in preceding chapters. Moreover, in developing this variant
of the perspective system the inevitable distortions were consciously shifted onto
the depiction of heights on the near and middle planes. If they were shifted
differently---correctly rendering heights---a different variant of the rigid
perspective system would be obtained, better suited for depicting tall architectural
structures on the near and middle planes. If, conversely, one requires that the
straight lines of objective space map to straight lines in the picture, and
distributes the inevitable distortions accordingly, the usual system of linear
perspective arises---which, incidentally, confirms from yet another standpoint the
thesis that this is merely a special case of the rigid perceptual perspective
system (though the table shows it to be far from the best way of distributing
distortions). With different solutions to the question of distributing the inevitable
distortions, other variants of rigid perspective systems can be obtained. But our
aim was the comparatively modest one of illustrating the fruitfulness of the
developed approach, not the creation of the various admissible geometries of visual
perception of perspective systems. (Appendix~8 gives formulas for variants of such
systems.)

\srcpage{231}

In the discussion of two systems of scientific perspective in Chapter~III, it was
noted that compensation for the distortions inherent in the linear perspective system
(far from complete) is achieved by the emphatic depiction of depth cues, including
non-geometric ones (aerial perspective, and so on). All of this can be employed in
the new perspective system considered here as well. One should bear in mind,
however, that a different geometry of depiction may also require a somewhat different
use of such non-geometric depth cues.

To conclude the comparison of the two systems of scientific perspective, let us note
their genetic kinship---both are central projections proceeding from the existence
of a single fixed viewpoint. They can therefore also be brought closer together in
nomenclature, by proposing the terms: system of central linear perspective and
system of central curvilinear perspective.

Let us draw some conclusions. First of all it should be emphasised once more that
the examination of Cézanne's paintings conducted above was limited to a single, and
possibly secondary, problem---the question of certain peculiarities of perspective
constructions. Cézanne himself placed first, following Delacroix, the ``composition
of colour'', and therefore an engagement with the geometric aspect of his works
alone, divorced from colour, can be regarded only as a preparatory stage that will
help toward a deeper study of his art. Moreover, even the geometric questions were
considered in a narrowly restricted manner; in particular, his thesis that in nature
all objects are modelled from the sphere, the cone, and the cylinder---a thesis
connected with his views on the volumetric modelling of forms---was not discussed
at all.

The comparison of the general principles of Cézanne's perspective constructions with
the perceptual perspective system of appropriate type, conducted above, showed their
striking closeness. What has been said allows the prevalent view of this side of
Cézanne's work to be changed---a view formulated by Novotny as follows:
``The perspective system that collectively weakened the rendering of spatiality given
by exact linear perspective construction is hardly a perspective at all, and therefore
it is appropriate to ask whether the end of scientific perspective is already
observable in Cézanne's painting rather than later, in the conspicuous forms of
abstract painting.''\footnote{F.~Novotny. Op.~cit., S.\,144.}

\srcpage{232}

It is now clear that associating Cézanne with the origins of abstract painting---at
any rate on the basis of his perspective constructions---is completely unjustified.
Rather the opposite should be asserted: Cézanne took the next step in perspective
construction methods. He moved from the use of the ordinary system of linear
perspective to the fuller and in certain respects more perfect system of perceptual
perspective, no less scientific than the one from whose abandonment the artist sought
to strengthen the expressive power of his paintings. Therefore Cézanne's name should
be associated not with the \emph{destruction} of ``scientific perspective'' but with
the \emph{development} of scientific perspective---even if one accomplished by the
artist intuitively, without knowledge of the sufficiently complex mathematical
foundations underlying the perspective constructions he employed.

If one likes, Cézanne's painting proved to be more realistic---more closely
corresponding to natural visual perception---than the canvases of artists who
adhered to the classical system of linear perspective.
